\documentclass[11pt,a4paper]{article}
\usepackage{notesthemejambro}
\renewcommand{\memouser}{ACB}      % default username (your initials)
\renewcommand{\memocolor}{BoEred}  % default color

%==================================================
\title{\vspace{-2cm}\color{BoEblue}
\textbf{notesthemejambro.sty}\\[0.3em]
\Large Feature Demo\\[0.3em]
\large v1.1-dev --- \textit{\today}
}
\author{}
\date{}
\begin{document}
\maketitle
\tableofcontents
\newpage
\vspace{-1cm}
%==================================================

%==================================================
\section{Typography and Fonts}
%==================================================

\subsection{Body Text}
%--------------------------------------------------

The default font is \textbf{Cabin} (sans-serif), loaded via \verb|[sfdefault]{cabin}|. All body text, headings, and captions use this family. Math is rendered in the standard Computer Modern math fonts alongside Cabin.

A short paragraph: The quick brown fox jumps over the lazy dog. Numbers: 0123456789. Greek: $\alpha, \beta, \gamma, \delta, \epsilon$. Operators: $\sum_{i=1}^{n} x_i$, $\int_0^\infty e^{-x}\,dx$.

\subsection{Handwriting Font}
%--------------------------------------------------

The Augie font is available for hand-drawn effects:

\begin{itemize}
    \item {\hand This sentence is set in the Augie handwriting font.}
    \item \handbold{This sentence uses handbold — bold Augie via PDF stroke literal.}
    \item Mixed: regular text alongside {\hand handwritten words} in the same line.
\end{itemize}

\subsection{Pencil Underlines}
%--------------------------------------------------

\lapis{Default underline} uses the \texttt{carandache} color (goldenrod). A custom color is passed as an optional argument:

\begin{itemize}
    \item \lapis{carandache gold} — default
    \item \lapis[tomato]{tomato red}
    \item \lapis[note]{note blue}
    \item \lapis[cobalt]{cobalt teal}
    \item \lapis[BoEpurple]{BoE purple}
\end{itemize}

Underlines work inside longer sentences too: the model predicts that \lapis{output falls} whenever the borrowing constraint \lapis[tomato]{binds}.

%==================================================
\section{Pencil Boxes}
%==================================================

\subsection{Inline Boxes}
%--------------------------------------------------

\lapisbox{Default gold box} uses no options. Custom colors:

\begin{itemize}
    \item \lapisbox{gold} — default
    \item \lapisbox[color=tomato]{tomato}
    \item \lapisbox[color=note]{note blue}
    \item \lapisbox[color=cobalt]{cobalt}
    \item \lapisbox[color=BoEpurple]{BoE purple}
\end{itemize}

Boxes work mid-sentence: the key result is that \lapisbox[color=note]{$\beta > 0$} whenever agents are impatient.

\subsection{Equation Boxes}
%--------------------------------------------------

\verb|\lapisboxeq| wraps display math in a pencil box without requiring the caller to enter math mode:
\[
    \lapisboxeq{Y_t = C_t + I_t + G_t + NX_t}
    \qquad
    \lapisboxeq[color=tomato]{r_t = \rho + \phi_\pi(\pi_t - \bar\pi) + \phi_y \tilde y_t}
\]

It also works inline: the steady-state condition is \lapisboxeq{r = \rho}.

\subsection{Labeled boxes and pencil arrows}
%--------------------------------------------------

Boxes with \verb|label=| create named TikZ nodes. A later \texttt{tikzpicture} can target them; compile \textbf{twice} for cross-picture references to resolve.


\begin{itemize}
    \item The \lapisbox[color=note,label=demand]{aggregate demand} externality 
    \item The \lapisbox[color=tomato,label=firesale]{firesale} externality
\end{itemize}
\begin{tikzpicture}[tpstyle]
    \draw[pencil,->,tomato!80] ([yshift=-2pt]firesale.south east) to[bend right] +(+0.7,-0.5) node[anchor=west,opacity=1,tomato] {\hand A comment};
\end{tikzpicture}

%==================================================
\section{Annotations and Notes}
%==================================================

\subsection{Inline Todo Notes}
%--------------------------------------------------

\NB{This is an inline note produced by the \texttt{NB} command. It uses the \texttt{note} blue background from \texttt{todonotes}.}

Todo notes are useful for flagging gaps during drafting. They sit in the text flow (not the margin) and are easy to search.

\subsection{Draft Annotations}
%--------------------------------------------------

\memo{Default user (ACB, set once in the preamble) and default color (BoEred).}

\memo[user=JB,color=note]{Custom user and color via optional key=value argument.}

\memo[color=cobalt]{Color-only override — username still comes from \texttt{\string\memouser}.}

Regular text continues after the annotation without extra spacing.

\subsection{Note Boxes}
%--------------------------------------------------

A \verb|notebox| is a breakable shaded callout:

\begin{notebox}
    \textbf{Key result.} In the Korinek--Simsek (2016) model, competitive borrowers do not internalise that their deleveraging reduces aggregate income. The social planner internalises this and chooses a lower debt level in period 0.
\end{notebox}

Note boxes accept arbitrary content including math and itemize:

\begin{notebox}
    \textbf{Two externalities.}
    \begin{itemize}
        \item \textbf{Aggregate demand}: deleveraging $\Rightarrow$ lower income $\Rightarrow$ tighter constraint.
        \item \textbf{Firesale}: forced selling $\Rightarrow$ lower asset prices $\Rightarrow$ lower collateral.
    \end{itemize}
    Both generate over-borrowing in equilibrium.
\end{notebox}

%==================================================
\section{Mathematics}
%==================================================

\subsection{Standard Environments}
%--------------------------------------------------

The \verb|\E| shorthand gives the expectations operator:
\[
    \E_t\left[\frac{u'(C_{t+1})}{u'(C_t)}\frac{1+r_{t+1}}{1+\pi_{t+1}}\right] = 1
\]

A multi-line derivation using \texttt{align}:
\begin{align}
    C_t^b + B_{t+1} &= y_t + (1+r_t)B_t \\
    B_{t+1}         &\leq \phi \label{eq:constraint}\\
    \lambda_t       &= u'(C_t^b) - \beta^b(1+r_t)\E_t[u'(C_{t+1}^b)] \geq 0
\end{align}

%==================================================
\section{Tables}
%==================================================

\subsection{Standard Tabularx}
%--------------------------------------------------

The \texttt{Y} column type is a centered auto-width \texttt{X} column. The style deliberately redefines \verb|\notesize| to produce compact 8.5pt annotation text below the table.

\begin{table}[h!]
\centering
\caption{Summary of externality types}
\begin{tabularx}{0.85\linewidth}{lYY}
    \toprule
    & \textbf{Aggregate demand} & \textbf{Firesale} \\
    \midrule
    Mechanism   & Spending $\to$ income & Prices $\to$ collateral \\
    Affected by & Borrower deleveraging & Asset sales \\
    Policy tool & Debt tax / limit & LTV cap \\
    Model       & Korinek--Simsek & Lorenzoni \\
    \bottomrule
\end{tabularx}
\\[0.4em]
{\notesize Note: Both mechanisms generate over-borrowing in competitive equilibrium.}
\end{table}

\subsection{Table with Row Colors}
%--------------------------------------------------

The \texttt{[table]} option on \texttt{xcolor} enables \verb|\rowcolor|:

\begin{table}[h!]
\centering
\caption{Model parameters}
\begin{tabularx}{0.7\linewidth}{lYY}
    \toprule
    \textbf{Parameter} & \textbf{Symbol} & \textbf{Value} \\
    \midrule
    \rowcolor{note!8} Discount factor (borrower) & $\beta^b$ & 0.90 \\
    Discount factor (lender)  & $\beta^l$ & 0.95 \\
    \rowcolor{note!8} Borrowing limit           & $\phi$    & 0.75 \\
    Frisch elasticity         & $\varphi$ & 1.00 \\
    \rowcolor{note!8} Price rigidity            & $\kappa$  & 0.50 \\
    \bottomrule
\end{tabularx}
\end{table}

%==================================================
\section{Paragraph Numbering}
%==================================================

Numbered paragraphs use \verb|\para|. The counter increments automatically and typesets a number followed by a quad of space.

\para This is the first numbered paragraph. The number is produced by \verb|\para| and counts independently of section numbers. It is useful for notes that cross-reference specific arguments.

\para This is the second numbered paragraph. The borrowing constraint $B_{t+1} \leq \phi$ is assumed to bind in period 1 but not in period 0.

\para This is the third numbered paragraph. The social planner internalises the aggregate demand externality and imposes a tax $\tau^*$ on period-0 borrowing.

\begin{notebox}
    Paragraphs~1--3 above illustrate that \verb|\para| works alongside \verb|notebox| and other environments. Cross-referencing via \verb|\label| / \verb|\ref| works because \verb|\para| calls \verb|\refstepcounter|.
\end{notebox}

%==================================================
\section{Pencil Arrows}
%==================================================

\subsection{Inline Directional Arrows}
%--------------------------------------------------

The \verb|\jarrow*| commands produce small baseline-aligned pencil arrows. All four directions are available:

\begin{itemize}
    \item Up:    GDP growth \jarrowup\ accelerated in Q3.
    \item Down:  Inflation \jarrowdown\ fell below target.
    \item Right: The effect propagates \jarrowright\ through the banking sector.
    \item Left:  Demand shifted \jarrowleft\ after the shock.
\end{itemize}

Custom options via key=value: \jarrowup[color=tomato]\ (tomato), \jarrowr[len=0.6cm,lw=0.1cm]\ (long, thick), \jarrowd[color=note,pad=0.1em]\ (tight blue).

\subsection{Overlay Arrows with Marker Nodes}
%--------------------------------------------------

\verb|\marker{name}{text}| wraps visible text in a named TikZ node. A subsequent overlay \texttt{tikzpicture} can draw an \texttt{arrow}-style line to it. Compile \textbf{twice} for cross-picture references to resolve.

Consider the \marker{fisherA}{Fisher equation} and the \marker{quantA}{quantity theory}:
\begin{tikzpicture}[tpstyle]
    \draw[arrow] ([yshift=-3pt]fisherA.south) to[bend right=20] +(+0.2,-0.6)
        node[anchor=north west] {\hand links $r$ and $\pi$};
    \draw[arrow,note] ([yshift=-3pt]quantA.south) to[bend right=20] +(-0.2,-0.6)
        node[anchor=north west] {\hand links $M$ and $P$};
\end{tikzpicture}
\vspace{1.5cm}

%==================================================
\section{Package Options}
%==================================================

\subsection{Wiggle}
%--------------------------------------------------

The pencil amplitude is controlled by the \verb|wiggle=| package option:

\begin{itemize}
    \item \verb|\usepackage{notesthemejambro}| $\to$ default amplitude (0.4pt)
    \item \verb|\usepackage[wiggle=0.3]{notesthemejambro}| $\to$ subtle
    \item \verb|\usepackage[wiggle=1.5]{notesthemejambro}| $\to$ expressive
\end{itemize}

This document uses the default. Recompile with a different option to compare:
\[
    \lapisboxeq{Y = \alpha + \beta X + \varepsilon} \qquad
    \lapis{underline at default wiggle}
\]

\subsection{Memo Suppression}
%--------------------------------------------------

Use \verb|\usepackage[hidememo]{notesthemejambro}| when you want the same notes document without draft annotations. In that mode every \verb|\memo| call expands to nothing, while the rest of the style remains unchanged.

\end{document}
